% ===== PRÉAMBULE CANONIQUE — à coller VERBATIM en tête de CHAQUE .tex =====
\documentclass[11pt,a4paper]{article}
\usepackage[utf8]{inputenc}\usepackage[T1]{fontenc}\usepackage{lmodern}
\usepackage[french]{babel}
\usepackage{amsmath,amssymb,mathtools}\usepackage{icomma}
\usepackage[version=4]{mhchem}          % \ce{...} : équations & formules
\usepackage{siunitx}\sisetup{locale=FR} % \qty{}{}, \si{}, \ang{}
\usepackage{chemfig}                    % \chemfig{...} : structures développées
\usepackage{xcolor}\usepackage{colortbl}
\usepackage{tikz}\usepackage{pgfplots}\pgfplotsset{compat=1.18}
\usepackage[most]{tcolorbox}\usepackage{enumitem}\usepackage{array,booktabs}
\usepackage[a4paper,margin=2.2cm]{geometry}
\usetikzlibrary{arrows.meta,calc,angles,quotes,positioning,patterns,backgrounds,shapes.geometric}
\definecolor{primary}{RGB}{222,49,99}\definecolor{primarydark}{RGB}{150,22,55}
\definecolor{defcol}{RGB}{0,114,178}\definecolor{methcol}{RGB}{52,92,148}
\definecolor{excol}{RGB}{0,135,90}\definecolor{attncol}{RGB}{200,40,55}
\tcbset{boxbase/.style={enhanced,breakable,boxrule=0.9pt,arc=2.5pt,
  left=3mm,right=3mm,top=2.3mm,bottom=2.3mm,fonttitle=\bfseries,coltitle=white}}
% --- boîtes TUTORIEL (noms EXACTS, ne pas renommer) ---
\newtcolorbox{cle}[1][]{boxbase,colback=defcol!6,colframe=defcol,
  title={\textsc{À retenir}\ifx\relax#1\relax\relax\else\textmd{ : \itshape #1}\fi}}
\newtcolorbox{methode}[1][]{boxbase,colback=methcol!7,colframe=methcol,sharp corners,
  title={\textsc{Méthode}\ifx\relax#1\relax\relax\else\textmd{ --- \itshape #1}\fi}}
\newtcolorbox{exemplebox}[1][]{boxbase,colback=excol!5,colframe=excol,
  title={\textsc{Exemple}\ifx\relax#1\relax\relax\else\textmd{ : \itshape #1}\fi}}
\newtcolorbox{piege}[1][]{boxbase,colback=attncol!6,colframe=attncol,
  title={\textsc{Piège}\ifx\relax#1\relax\relax\else\textmd{ : \itshape #1}\fi}}
\newtcolorbox{remarque}[1][]{boxbase,colback=black!4,colframe=black!45,
  title={\textsc{Remarque}\ifx\relax#1\relax\relax\else\textmd{ : \itshape #1}\fi}}
% --- boîtes EXERCICE / CORRIGÉ (noms EXACTS, auto-comptées) ---
\newtcolorbox[auto counter]{exo}[1][]{boxbase,colback=primary!4,colframe=primary,
  title={\textsc{Exercice}~\thetcbcounter\ifx\relax#1\relax\relax\else\textmd{ --- \itshape #1}\fi}}
\newtcolorbox[auto counter]{corrige}[1][]{boxbase,colback=excol!4,colframe=excol,
  title={\textsc{Corrigé}~\thetcbcounter\ifx\relax#1\relax\relax\else\textmd{ --- \itshape #1}\fi}}
\newcommand{\R}{\mathbb{R}}\newcommand{\N}{\mathbb{N}}\newcommand{\Z}{\mathbb{Z}}\newcommand{\ds}{\displaystyle}
% (un \usepackage{fancyhdr}+en-tête est possible ; il est ignoré côté web)
% =========================================================================
\begin{document}
\sisetup{per-mode=symbol}
\begin{center}\color{primarydark}\large\bfseries Thème 5 — Corrigé \quad$\bullet$\quad Niveau Moyen\end{center}

\begin{corrige}[décrire complètement une pile]
\begin{enumerate}[label=\alph*)]
  \item Cathode $=$ \ce{Cu} ($+0{,}34$, pôle $+$) ; anode $=$ \ce{Ni} ($-0{,}26$, pôle $-$).
  \item Anode : \ce{Ni -> Ni^{2+} + 2 e^-} ; cathode : \ce{Cu^{2+} + 2 e^- -> Cu}. Globale :
        \ce{Ni + Cu^{2+} -> Ni^{2+} + Cu}.
  \item \ce{(-) Ni | Ni^{2+} || Cu^{2+} | Cu (+)} ; $\Delta E^{\circ}=0{,}34-(-0{,}26)=+0{,}60$~V.
\end{enumerate}
\end{corrige}

\begin{corrige}[que devient chaque électrode ?]
\begin{enumerate}[label=\alph*)]
  \item Les anions \ce{NO3^-} migrent vers l'\textbf{anode} (compartiment \ce{Cu}) ; les cations
        \ce{K+} vers la \textbf{cathode} (compartiment \ce{Ag}). Cela maintient l'électroneutralité.
  \item L'argent est la \textbf{cathode} : \ce{Ag+ + e^- -> Ag}, l'argent se dépose, l'électrode
        \textbf{s'épaissit}.
  \item Le cuivre est l'\textbf{anode} : \ce{Cu -> Cu^{2+} + 2 e^-}, le métal se dissout, l'électrode
        \textbf{s'amincit}.
\end{enumerate}
\end{corrige}

\begin{corrige}[masse déposée par électrolyse]
$t=\qty{1}{h}=\qty{3600}{s}$, $n=3$ :
\[
m=\frac{M\,I\,t}{n\,F}=\frac{27\times 5{,}0\times 3600}{3\times 96485}
=\frac{486\,000}{289\,455}\approx\qty{1.68}{g}.
\]
\end{corrige}

\begin{corrige}[combien de temps ?]
On isole $t$ : $t=\dfrac{m\,n\,F}{M\,I}=\dfrac{10\times 1\times 96485}{108\times 5{,}0}
=\dfrac{964\,850}{540}\approx\qty{1787}{s}\approx\qty{29.8}{min}$ (environ une demi-heure).
\end{corrige}

\begin{corrige}[pile ou électrolyse ?]
\begin{enumerate}[label=\alph*)]
  \item Une \textbf{pile} exploite une réaction \textbf{spontanée} ($\Delta E^{\circ}>0$) pour produire
        du courant ; une \textbf{électrolyse} \textbf{force} une réaction \textbf{non spontanée} grâce
        à un \textbf{générateur} qui impose le courant.
  \item La réduction a lieu à la \textbf{cathode}, reliée au pôle \textbf{$-$} du générateur.
  \item L'oxydation a lieu à l'\textbf{anode} (reliée au pôle $+$).
\end{enumerate}
\end{corrige}

\end{document}
